\section{Cualitativas}
En esta sección se abordarán las variables cualitativas, que son aquellas que describen
características o atributos de los elementos de estudio. Estas variables no tienen
un valor numérico intrínseco y se dividen en dos subcategorías principales: nominales y ordinales.
Cómo por sí solas no permiten realizar un análisis estadístico profundo, nos limitaremos a contar la
cantidad de elementos que pertenecen a cada categoría y encontrar la moda, que es la categoría con
mayor frecuencia.

\subsection{Nominales}
En el análisis vamos a trabajar con \Variable{'gender'}.

\subsubsection{Frecuencias de la variable gender}
\begin{DatosTabla}{L{0.20\linewidth} C{0.14\linewidth} C{0.14\linewidth} C{0.14\linewidth} C{0.14\linewidth}}{\textbf{Variable} & \textbf{fi} & \textbf{fr} & \textbf{Fi} & \textbf{Fr}}

\TablaClave{male} & 53 & 47,75\% & 53 & 47,75\% \\
\TablaClave{female} & 58 & 52,25\% & 111 & 100,00\% \\
\textbf{Total} & \textbf{111} & \textbf{100,00\%} & \textbf{111} & \textbf{100,00\%} \\

\end{DatosTabla}


\AnalisisDescriptivo
\PalabraClave{Moda}: female (58 apariciones, 52,3\%) 

\TituloGrafico{male vs female}
% =====================================================================
%  figures/gender-torta.tex
%  Gráfico de torta para la variable 'gender'.
%  (Reemplaza a img/grafico_gender.png)
% =====================================================================
\begin{figure}[H]
\centering
\begin{tikzpicture}[scale=0.9]
  \pie[
    text=legend,
    radius=2.6,
    color={azul, naranja},
    sum=auto,
    after number=\%,
  ]{
    52.25/female,
    47.75/male
  }
\end{tikzpicture}
\caption{Gráfico de torta para la variable \Variable{gender}.}
\label{fig:grafico_gender}
\end{figure}


% --------------------------------------------------------------------------------------------------

\subsection{Ordinales}
Ahora vamos a trabajar con \Variable{'parent\_max\_education'}.

\subsubsection{Frecuencias de la variable parent\_max\_education}
\begin{DatosTabla}{L{0.20\linewidth} C{0.14\linewidth} C{0.14\linewidth} C{0.14\linewidth} C{0.14\linewidth}}{\textbf{Variable} & \textbf{fi} & \textbf{fr} & \textbf{Fi} & \textbf{Fr}}

\TablaClave{bachelor\_degree} & 104 & 93,69\% & 104 & 93,69\% \\
\TablaClave{master\_degree} & 7 & 6,31\% & 111 & 100,00\% \\
\textbf{Total} & \textbf{111} & \textbf{100,00\%} & \textbf{111} & \textbf{100,00\%} \\

\end{DatosTabla}


\AnalisisDescriptivo
\PalabraClave{Moda}: bachelor\_degree (104 apariciones, 93,69\%)

\TituloGrafico{Nivel educativo}
% =====================================================================
%  figures/nivel-educativo.tex
%  Gráfico de barras para 'parent_max_education'.
%  (Reemplaza a img/grafico_nivel_educativo.png)
% =====================================================================
\begin{figure}[H]
\centering
\begin{tikzpicture}
\begin{axis}[
  barras,
  every node near coord/.append style={font=\large},
  bar width=48pt,
  title={Nivel educativo máximo de los padres},
  ylabel={Cantidad de estudiantes},
  symbolic x coords={bachelor\_degree, master\_degree},
  xtick=data,
  ymin=0, ymax=120,
  enlarge x limits=0.5,
  nodes near coords,
]
\addplot coordinates {(bachelor\_degree,104) (master\_degree,7)};
\end{axis}
\end{tikzpicture}
\caption{Gráfico de barras para la variable \Variable{parent\_max\_education}.}
\label{fig:grafico_parent_max_education}
\end{figure}

